\documentclass[11pt]{article}
% Compile with pdfLaTeX from the repository root or from this folder.
\makeatletter
\def\input@path{{../}}
\makeatother
\usepackage{notes}
\setboolean{showanswer}{true}
\renewcommand{\solutionheight}{22mm}

\begin{document}
\lesson[Stewart, Calculus, 9th edition\enspace\textbullet\enspace Section 12.5]
  {1.5}{Equations of Lines}

\section{A point and a direction}
To describe a line in space, specify a point on the line and a nonzero
vector parallel to it. Starting at that point, move any scalar multiple of
the direction vector.

\begin{definition}[Vector equation of a line]
The line through $P_0=(x_0,y_0,z_0)$ with direction vector
$\B{v}=\cvec{a,b,c}\ne\B{0}$ has vector equation
\[
 \B{r}(t)=\B{r}_0+t\B{v},\qquad -\infty<t<\infty,
 \quad\text{where }\B{r}_0=\cvec{x_0,y_0,z_0}.
\]
The parameter $t$ labels points along the line; $t=0$ gives $P_0$.
\end{definition}

\begin{center}
\begin{tikzpicture}[scale=0.8,>=Latex,font=\small]
 \coordinate (O) at (0,0);
 \coordinate (P) at (1,2);
 \coordinate (Q) at (5,3);
 \draw[<->,notesrule,thick] (0,1.75)--(6,3.25) node[right,notesmuted]{$L$};
 \draw[->,thick,notesblue] (O)--(P) node[midway,above left]{$\B{r}_0$};
 \draw[->,very thick,notesteal] (P)--(Q) node[midway,above left]{$t\B{v}$};
 \draw[->,thick,notesink] (O)--(Q) node[midway,below right]{$\B{r}(t)$};
 \fill[notesink] (O) circle (1.7pt) node[below]{$O$};
 \fill[notesblue] (P) circle (2pt) node[below right]{$P_0$};
 \fill[notesteal] (Q) circle (2pt) node[above]{$P$};
\end{tikzpicture}
\end{center}

\subsection*{Parametric equations}
Reading the vector equation one component at a time gives
\[
 x=x_0+at,\qquad y=y_0+bt,\qquad z=z_0+ct.
\]
All three coordinates use the \textbf{same parameter value}. The parameter
need not represent time.

\begin{example}{Find vector and parametric equations for the line through
$P_0=(1,-2,3)$ parallel to $\B{v}=\cvec{2,1,-1}$.}
Use the point for $\B{r}_0$ and the given vector for the direction:
\[
 \B{r}(t)=\cvec{1,-2,3}+t\cvec{2,1,-1}.
\]
The parametric equations are $x=1+2t$, $y=-2+t$, and $z=3-t$.
At $t=1$, the line passes through $(3,-1,2)$.
\end{example}

\subsection*{Symmetric equations}
When $a$, $b$, and $c$ are all nonzero, solve each coordinate equation for
$t$ and equate the results:
\[
 \frac{x-x_0}{a}=\frac{y-y_0}{b}=\frac{z-z_0}{c}.
\]
For Example 1, this becomes $(x-1)/2=y+2=(z-3)/(-1)$.

\newpage
\section{Building and interpreting line equations}
\subsection*{A line through two points}
For distinct points $A$ and $B$, the displacement $\vec{AB}$ supplies a
direction vector. Any nonzero scalar multiple of that vector can be used
instead.

\begin{example}{Find the line through $A=(1,-2,3)$ and $B=(5,0,-1)$.
How would you describe only the segment joining them?}
The direction vector is $\vec{AB}=\cvec{4,2,-4}$, giving
\[
 \B{r}(t)=\cvec{1,-2,3}+t\cvec{4,2,-4},\qquad t\in\mathbb{R}.
\]
Equivalently, the line has symmetric equations
\[
 \frac{x-1}{4}=\frac{y+2}{2}=\frac{z-3}{-4}.
\]
Restricting the parameter to $0\leq t\leq1$ gives the \textbf{line segment}:
$t=0$ gives $A$, $t=1$ gives $B$, and $t=1/2$ gives the midpoint $(3,-1,1)$.
\end{example}

\begin{definition}[Line segment]
If $\B{a}$ and $\B{b}$ are the position vectors of endpoints $A$ and $B$,
the segment from $A$ to $B$ is
\[
 \B{r}(t)=(1-t)\B{a}+t\B{b},\qquad 0\leq t\leq1.
\]
This is the same as $\B{r}(t)=\B{a}+t(\B{b}-\B{a})$.
\end{definition}

\subsection*{A zero direction component}
A zero component means that coordinate stays constant. Keep its equation
separate when writing symmetric equations; never divide by zero.

\begin{example}{A line passes through $(2,-1,4)$ with direction
$\cvec{0,3,-2}$. Write its equations and decide whether
$Q=(2,5,0)$ and $R=(2,5,2)$ lie on it.}
The parametric equations are
\[
 x=2,\qquad y=-1+3t,\qquad z=4-2t.
\]
A symmetric description is $x=2$, together with
\[
 \frac{y+1}{3}=\frac{z-4}{-2}.
\]
For $Q$, the $y$-coordinate gives $t=2$, which also gives $z=0$.
Thus $Q$ lies on the line.

For $R$, the $y$-coordinate still requires $t=2$, but then $z=0\ne2$.
Thus $R$ does not lie on the line. Each coordinate must agree at one
common parameter value.
\end{example}

\newpage
\section{How two lines can meet}
In a plane, two distinct lines are parallel or intersect. Space allows
one more possibility: lines that are neither parallel nor intersecting.

\begin{definition}[Relative positions of lines]
Two lines are \textbf{parallel} when their direction vectors are scalar
multiples of one another. Parallel lines may be distinct or coincident
(the same line).

Two nonparallel lines either \textbf{intersect} at one point or are
\textbf{skew}: they do not intersect and do not lie in a common plane.
\end{definition}

To test for an intersection, give each line its own parameter and equate
corresponding coordinates. A point of intersection must satisfy all three
coordinate equations.

\begin{example}{Determine whether these lines intersect:
\[
 L_1:\ \B{r}(t)=\cvec{1,0,2}+t\cvec{1,2,-1},\qquad
 L_2:\ \B{s}(u)=\cvec{0,3,0}+u\cvec{2,-1,1}.
\]}
Equating coordinates gives
\[
 1+t=2u,\qquad 2t=3-u,\qquad 2-t=u.
\]
Use the third equation in the second:
$2t=3-(2-t)=1+t$, so $t=1$ and $u=1$.
The first equation also holds: $1+1=2(1)$.

Substituting into either line gives the intersection point
$\boxed{(2,2,1)}$. Here the parameters happen to be equal; that need not
be true for other parametrizations.
\end{example}

\begin{example}{Classify the lines
$L_1:\ \B{r}(t)=\cvec{t,t,0}$ and
$L_2:\ \B{s}(u)=\cvec{u,-u,1}$.}
Their direction vectors $\cvec{1,1,0}$ and $\cvec{1,-1,0}$ are not scalar
multiples, so the lines are not parallel.

Every point of $L_1$ has $z=0$, while every point of $L_2$ has $z=1$.
The lines cannot intersect, so they are \textbf{skew}.
\end{example}

\subsection*{Different equations can describe the same line}
Changing the base point to another point on the line, or multiplying the
direction vector by a nonzero scalar, changes the parametrization without
changing the line. If two direction vectors are parallel, test whether a
point from one line lies on the other to distinguish coincident lines
from distinct parallel lines.
\end{document}
